\documentclass[12pt]{article}
\usepackage[utf8]{inputenc}
\usepackage[T1]{fontenc}
\usepackage{hyperref}
\title{Huron-Wyandot Language Resources for K-12 Curriculum Development}
\author{Compiled by Claude Sonnet 5.5 (claude-sonnet-5-5) for Andr\'as Kornai}
\date{}
\begin{document}
\maketitle

\section{Introduction}
Huron-Wyandot (ISO 639-3 wya; Glottocode wyan1247) is an Iroquoian language of North America. It is critically endangered. Its last native speaker reportedly died in 1972. Revival work continues in Oklahoma (Wyandotte Nation) and Quebec (Wendat Nation, Wendake). Sources disagree on terminology and on whether Wyandot and Wendat are separate languages or varieties of one language. Many documentary sources are old archival materials, so their spellings differ from current orthographies.

\section{Primary Educational Resources}
Search results indicate that the Wyandotte Nation (Oklahoma) offers Wyandot classes for grades K--4 in Wyandotte Public Schools and in its ``Turtle-Tots'' preschool. They also indicate that the nation has built online self-study lessons, though I could not confirm where those lessons are hosted. The Wendat Nation in Wendake offers adult and children's classes at the village school.
\begin{itemize}
\item Wendat language site (Quebec): \url{https://languewendat.com/en/?p=691}
\item \emph{DSMW -- A Wandat (Wyandot) Language Workbook} (2019, sold commercially with a CD; not produced by the tribe): \url{https://www.onlc.ca/bookstore/p/dsmw-a-wandat-wyandot-language-workbook}
\end{itemize}

\section{Linguistic and Reference Materials}
\begin{itemize}
\item Wikipedia overview of the Wyandot language: \url{https://en.wikipedia.org/wiki/Wyandot_language}
\item Kopris (2001), SUNY Buffalo dissertation with a grammar and a Huron-Wyandot--English wordlist, catalogued as \emph{A Grammar and Dictionary of Wyandot}: \url{https://wiki.ubc.ca/Documentation:RelLex/A_Grammar_and_Dictionary_of_Wyandot}
\item Steckley (2007), \emph{A Huron-English/English-Huron Dictionary}, based mainly on Potier's 1920 report: \url{https://wiki.ubc.ca/Documentation:RelLex/A_Huron-English/English-Huron_Dictionary}
\item Brebeuf-era French-Huron dictionary, published by Steckley in 2010: \url{https://knowledgebase.arts.ubc.ca/?p=3902}
\item Wendat Dictionary (UBC RelLex catalogue entry): \url{https://wiki.ubc.ca/Documentation:RelLex/Wendat_Dictionary}
\item Omniglot page on Wendat (Wa\textsuperscript{n}dat): \url{https://www.omniglot.com/writing/wyandot.htm}
\end{itemize}

\section{Religious and Community Resources}
\begin{itemize}
\item Wyandotte Nation culture page: \url{https://wyandotte-nation.org/culture/}
\item Profile of tribal linguist John Steckley, who teaches language courses to Wyandot youth (NCCIE): \url{https://www.nccie.ca/story/john-steckley-wyandot-tribal-linguist/}
\end{itemize}
I saw no scripture or hymn resources in the search results.

\section{Digital and Audio Resources}
\begin{itemize}
\item FirstVoices Huron-Wendat: an open-access multimedia wordlist with audio, images and text, covering Huron-Wyandot to and from English and French: \url{https://wiki.ubc.ca/Documentation:RelLex/Huron-Wendat_FirstVoices}
\end{itemize}

\section{Curriculum Development Context}
Classroom programmes exist at the Wyandotte Nation (K--4 and preschool). The tribe is reportedly developing a curriculum that would extend the language to elementary and high school. Wendake also has school classes. No complete published K--12 curriculum was found. The older dictionaries are useful as sources of vocabulary but need orthographic normalisation and review by the communities. Terminology and dialect differences between Oklahoma Wyandot and Quebec Wendat should be settled with the communities before materials are written.

\end{document}